\documentclass{article}
\usepackage[utf8]{inputenc}
\usepackage{amsmath}
\usepackage{geometry}
\geometry{margin=2cm}
\begin{document}

\section*{Inspeção de Bagagens por Raios X e a Remoção de Computadores Portáteis}

\subsection*{Interpretação do Enunciado e Estratégia}

O enunciado solicita explicar o motivo físico pelo qual passageiros devem retirar seus computadores portáteis (notebooks) das malas para passá-los isoladamente em inspeções por aparelhos de raios X em aeroportos.

\textbf{Termos importantes:}
\begin{itemize}
\item \textit{Explicação física:} Justificar o procedimento com base em princípios físicos da radiação e da interação com os materiais.
\item \textit{Esteira de inspeção:} Equipamento onde a bagagem é transportada para passar embaixo de aparelhos de raios X para análise.
\end{itemize}

A questão exige compreensão sobre como a radiação interage com materiais de diferentes densidades e a importância da separação para uma inspeção segura.

\textbf{Estratégia:}

Entender que aparelhos de raios X geram imagens baseadas na passagem da radiação que é bloqueada ou atenuada segundo a densidade dos objetos. Como computadores possuem componentes densos (metais e circuitos), eles criam sombras que impedem ver objetos atrás deles. É preciso identificar a alternativa que explicita essa razão para a remoção dos computadores.

\subsection*{Conteúdos Relevantes e Mini Revisão}

\begin{tabular}{p{5cm} p{6cm} p{5cm}}
\textbf{Conceito} & \textbf{Ideia-chave} & \textbf{Aplicação} \\
\hline
Interação de raios X com a matéria & Materiais densos bloqueiam mais a radiação & Componentes metálicos bloqueiam a passagem dos raios X \\
Formação de imagens por raios X & Contraste da passagem da radiação gera a imagem & Objetos densos aparecem mais na imagem, podendo ocultar itens atrás \\
Procedimento em aeroportos & Separar itens para melhor inspeção & Retirar computadores garante imagem clara para todos os objetos \\
\end{tabular}

\
\textbf{Teoria resumida:}

\begin{itemize}
\item Aparelhos de raios X em aeroportos emitem radiação capaz de atravessar os objetos da bagagem, gerando imagens conforme atenuação da radiação.
\item Materiais mais densos, como metais e circuitos eletrônicos, bloqueiam ou atenuam fortemente a radiação.
\item Com computadores dentro da bagagem, eles criam sombras impendindo a visualização clara de objetos atrás.
\item Removê-los permite inspeção isolada e completa de todos itens.
\end{itemize}

\subsection*{Resolução e \textit{Pulo do Gato}}

Para responder:
\begin{enumerate}
\item Reconheça que raios X geram imagens por diferenças na densidade dos objetos.
\item Identifique que componentes do computador têm alta densidade e bloqueiam os raios X.
\item Compreenda que tal bloqueio ocultaria objetos que estejam atrás do computador dentro da bagagem.
\item Conclua que, para garantir a inspeção completa, os computadores são retirados e inspecionados isoladamente.
\end{enumerate}

\textbf{Pulo do gato:} Itens muito densos bloqueiam a radiação dos raios X e criam sombras que podem esconder objetos; portanto, para uma inspeção eficaz, esses itens devem ser removidos para exame separado.

\subsection*{Análise das Alternativas e Armadilhas}

\begin{itemize}
\item \textbf{Alternativa A:} Errada, pois metais interagem com raios X e aparecem na imagem, bloqueando a radiação.
\item \textbf{Alternativa B:} Incorreta, pois raios X não causam desmagnetização nos discos rígidos dos computadores.
\item \textbf{Alternativa C:} Errônea, porque raios X não têm capacidade significativa para aquecer ou refletir metais como luz visível.
\item \textbf{Alternativa D:} Correta, por explicar que a alta densidade bloqueia raios X, prejudicando a visualização e justificando remoção.
\item \textbf{Alternativa E:} Errada, por atribuir efeitos incorretos de difração e ionização que não impedem a inspeção.
\end{itemize}

\subsection*{Código da Questão}

\texttt{CN\_FISICA\_ENSINO\_MEDIO\_001}

\end{document}